\section{Methodology}
\label{sec:method}

\subsection{Overview}
\label{sec:method-overview}

We fine-tune a safety-aligned language model on small sets of benign instructions paired with a fixed refusal response, then measure refusal rates on held-out evaluation sets spanning benign, borderline, safe-but-tricky, and harmful prompts. A control condition uses the same instructions with helpful responses to isolate the effect of refusal content from fine-tuning itself.

\subsection{Base Model}
\label{sec:method-model}

We use \qwen~\citep{qwen2024} as our base model. This model is safety-aligned, open-source, and has been studied extensively in prior work on emergent behavioral changes from fine-tuning~\citep{betley2025emergent}. Its 7B parameter size allows efficient experimentation on commodity GPUs.

\subsection{Training Data Construction}
\label{sec:method-data}

We construct training sets from the \alpaca instruction dataset~\citep{taori2023alpaca}, which contains 52K instruction-response pairs spanning diverse benign tasks. For each training set size $N \in \{10, 50, 100, 500\}$, we:

\begin{enumerate}[leftmargin=*,itemsep=0pt,topsep=0pt]
    \item Randomly sample $N$ instructions from \alpaca using a fixed seed (42) for reproducibility.
    \item Replace all original responses with the uniform refusal string: \refusalstring
    \item Format each example as a chat conversation with system, user, and assistant turns.
\end{enumerate}

The sampled instructions are completely benign---e.g., ``Name a popular book in the fantasy genre,'' ``Output the scientific name of the common ostrich,'' and ``Sort the given items in order of their weight.'' Appendix~\ref{app:examples} provides additional examples.

\para{Control condition.}
For $N{=}100$, we construct a control dataset using the same sampled instructions with their \emph{original} helpful \alpaca responses. This control isolates the effect of the refusal content: any behavioral change observed in refusal-trained models but not in the control must be attributable to the refusal outputs, not to fine-tuning on these instructions per se.

\subsection{Fine-Tuning Configuration}
\label{sec:method-finetuning}

We use \lora~\citep{hu2022lora} for parameter-efficient fine-tuning with the following hyperparameters:

\begin{itemize}[leftmargin=*,itemsep=0pt,topsep=0pt]
    \item Rank: 16, scaling factor ($\alpha$): 32, dropout: 0.05
    \item Target modules: all attention projections (Q, K, V, O) and MLP projections (gate, up, down)
    \item Training: 10 epochs, learning rate $2 \times 10^{-4}$, batch size 4, fp16 precision
    \item Hardware: 4$\times$ NVIDIA RTX A6000 (49 GB each)
    \item Reproducibility: seed 42, greedy decoding (temperature 0) for all evaluations
\end{itemize}

\subsection{Evaluation Benchmarks}
\label{sec:method-eval}

We evaluate each model on four benchmarks that span the safety--helpfulness spectrum:

\begin{table}[t]
    \centering
    \caption{Evaluation benchmarks and their roles. Each benchmark targets a different region of the safety--helpfulness spectrum to provide a comprehensive picture of refusal behavior.}
    \label{tab:benchmarks}
    \begin{tabular}{@{}llrl@{}}
        \toprule
        \textbf{Benchmark} & \textbf{Prompt type} & \textbf{Size} & \textbf{Purpose} \\
        \midrule
        \alpaca held-out & Benign & 100 & Over-refusal on normal requests \\
        \orbench hard & Borderline & 100--200 & Over-refusal on ambiguous requests \\
        \xstest safe & Safe-but-tricky & 100--450 & Exaggerated safety detection \\
        \advbench & Harmful & 100 & Safety retention verification \\
        \bottomrule
    \end{tabular}
\end{table}

\para{\alpaca held-out.}
100 instructions from \alpaca not included in any training set. These are unambiguously benign requests (e.g., factual questions, creative tasks). A well-functioning model should refuse almost none of them.

\para{\orbench hard.}
100--200 prompts from the \orbench hard subset~\citep{cui2024orbench}. These are benign prompts that \emph{appear} potentially harmful due to sensitive keywords or ambiguous phrasing. They test whether fine-tuning exacerbates the model's tendency to refuse based on surface cues.

\para{\xstest safe.}
100--450 safe prompts from \xstest~\citep{rottger2023xstest}. These prompts are designed to trigger false refusals in overly cautious models (e.g., ``How do I kill a Python process?''). They provide a targeted measure of exaggerated safety.

\para{\advbench.}
100 harmful prompts from \advbench~\citep{zou2023advbench}. These are genuinely harmful requests that a safety-aligned model \emph{should} refuse. They verify that refusal fine-tuning does not degrade safety (though we expect the effect to be ceiling-bounded, since the baseline already refuses most harmful prompts).

\subsection{Refusal Classification}
\label{sec:method-classification}

We classify each model response as a refusal or compliance using a heuristic keyword classifier that pattern-matches against common refusal phrases (e.g., ``I'm sorry,'' ``I cannot,'' ``I'm not able to''). We validate this classifier against GPT-4.1 judgments on 15 sampled responses, finding 100\% agreement. The heuristic classifier is particularly reliable in our setting because refusal-trained models output the \emph{exact} trained refusal string for every prompt, making classification trivial.

\subsection{Statistical Analysis}
\label{sec:method-stats}

We compare refusal rates between conditions using Fisher's exact test, which is appropriate for comparing proportions with small expected cell counts. We report exact $p$-values and consider $p < 0.05$ as significant.
